% GENERATED FILE. Edit canonical lessons, then run npm run generate:books.
% canonical-source-hash: fnv1a64:049eab6b4c504952
% canonical-lessons: BN-C545-onyotom, BN-C545-ogradhikar, BN-C545-sadhyo, BN-C545-mitobyoyi, BN-C545-thaorano, BN-R545-review-539-545, BN-R545-review-final-545

\chapter{Choices}
\label{ch:bn-choices}
\hlchaptermodality{545}

\begin{chapteropening}
\textbf{By the end of this chapter:} \emph{I can talk about priorities, thrift and what is within my power.}

Complete the last lesson of chapter 545: I can talk about priorities, thrift and what is within my power.
\end{chapteropening}

\section[\texorpdfstring{\bn{অন্যতম} (ônyôtôm)}{অন্যতম (ônyôtôm)}]{\bn{অন্যতম} (ônyôtôm) — one of several}

\label{lesson:BN-C545-onyotom}



\begin{quote}
Before the new one: say the Bengali for a scramble, then the Bengali for in a jiffy.
\end{quote}

\subsection*{You'll want to know: \bn{অন্যতম}}
\textbf{\bn{অন্যতম}} — \emph{ônyôtôm} — \textquotedblleft{}one of several\textquotedblright{}.

\begin{quote}
\bn{ইলিশ} \bn{আমার} \bn{প্রিয়} \bn{মাছগুলোর} \bn{অন্যতম।}
\end{quote}

\begin{quote}
\emph{Ilish āmār priyô māchhgulor ônyôtôm.}
\end{quote}

\begin{quote}
Hilsa is one of my favourite fish.
\end{quote}

\begin{cousinweb}
\textbf{\bn{অন্যতম}} (\emph{ônyôtôm}) is a tatsama, Sanskrit \emph{anyatama}, \textquotedblleft{}any one of many\textquotedblright{} (Monier-Williams).
\end{cousinweb}

\begin{grammarlens}[title={What you've built}]
One more word for saying what happens if.
\end{grammarlens}

\subsection*{Your turn}
\begin{itemize}
\raggedright
  \item \emph{Say it:} \emph{ônyôtôm}
  \item \emph{Say it:} \emph{ônyôtôm}, once more
  \item \emph{Say it:} \emph{Ilish āmār priyô māchhgulor ônyôtôm.} — Hilsa is one of my favourite fish.
  \item \emph{Recall:} say the Bengali for a scramble, then the Bengali for in a jiffy, then say \emph{ônyôtôm} again
\end{itemize}

\begin{tcolorbox}[breakable,colback=teal!4,colframe=teal!35!black,title={Before you move on}]
What does \bn{অন্যতম} mean? (\textquotedblleft{}One of several\textquotedblright{}.) Say it once more.
\end{tcolorbox}
\section[\texorpdfstring{\bn{অগ্রাধিকার} (ôgrādhikār)}{অগ্রাধিকার (ôgrādhikār)}]{\bn{অগ্রাধিকার} (ôgrādhikār) — priority}

\label{lesson:BN-C545-ogradhikar}



\begin{quote}
Before the new one: say the Bengali for in a jiffy, then the Bengali for one of several.
\end{quote}

\subsection*{You'll want to know: \bn{অগ্রাধিকার}}
\textbf{\bn{অগ্রাধিকার}} — \emph{ôgrādhikār} — \textquotedblleft{}priority\textquotedblright{}.

\begin{quote}
\bn{এই} \bn{হাসপাতালে} \bn{শিশুরা} \bn{অগ্রাধিকার} \bn{পায়।}
\end{quote}

\begin{quote}
\emph{Ei hāspātāle shishurā ôgrādhikār pāy.}
\end{quote}

\begin{quote}
In this hospital children get priority.
\end{quote}

\begin{cousinweb}
\textbf{\bn{অগ্রাধিকার}} (\emph{ôgrādhikār}) joins Sanskrit \emph{agra}, \textquotedblleft{}foremost\textquotedblright{}, and \emph{adhikāra}, \textquotedblleft{}authority, a claim\textquotedblright{} (Monier-Williams): the first claim.
\end{cousinweb}

\begin{grammarlens}[title={What you've built}]
One more word for saying what happens if.
\end{grammarlens}

\subsection*{Your turn}
\begin{itemize}
\raggedright
  \item \emph{Say it:} \emph{ôgrādhikār}
  \item \emph{Say it:} \emph{ôgrādhikār}, once more
  \item \emph{Say it:} \emph{Ei hāspātāle shishurā ôgrādhikār pāy.} — In this hospital children get priority.
  \item \emph{Recall:} say the Bengali for in a jiffy, then the Bengali for one of several, then say \emph{ôgrādhikār} again
\end{itemize}

\begin{tcolorbox}[breakable,colback=teal!4,colframe=teal!35!black,title={Before you move on}]
What does \bn{অগ্রাধিকার} mean? (\textquotedblleft{}Priority\textquotedblright{}.) Say it once more.
\end{tcolorbox}
\section[\texorpdfstring{\bn{সাধ্য} (sādhyô)}{সাধ্য (sādhyô)}]{\bn{সাধ্য} (sādhyô) — one's power; what one can manage}

\label{lesson:BN-C545-sadhyo}



\begin{quote}
Before the new one: say the Bengali for one of several, then the Bengali for priority.
\end{quote}

\subsection*{You'll want to know: \bn{সাধ্য}}
\textbf{\bn{সাধ্য}} — \emph{sādhyô} — \textquotedblleft{}one's power; what one can manage\textquotedblright{}.

\begin{quote}
\bn{আমার} \bn{সাধ্যমতো} \bn{আমি} \bn{তোমাকে} \bn{সাহায্য} \bn{করব।}
\end{quote}

\begin{quote}
\emph{Āmār sādhyômôto āmi tomāke sāhājyô kôrbô.}
\end{quote}

\begin{quote}
I'll help you as much as I can.
\end{quote}

With -\bn{মতো}: \bn{সাধ্যমতো}, as far as one can.

\begin{cousinweb}
\textbf{\bn{সাধ্য}} (\emph{sādhyô}) is a tatsama, Sanskrit \emph{sādhya}, \textquotedblleft{}to be accomplished, attainable\textquotedblright{} (Monier-Williams).
\end{cousinweb}

\begin{grammarlens}[title={What you've built}]
One more word for saying what happens if.
\end{grammarlens}

\subsection*{Your turn}
\begin{itemize}
\raggedright
  \item \emph{Say it:} \emph{sādhyô}
  \item \emph{Say it:} \emph{sādhyô}, once more
  \item \emph{Say it:} \emph{Āmār sādhyômôto āmi tomāke sāhājyô kôrbô.} — I'll help you as much as I can.
  \item \emph{Recall:} say the Bengali for one of several, then the Bengali for priority, then say \emph{sādhyô} again
\end{itemize}

\begin{tcolorbox}[breakable,colback=teal!4,colframe=teal!35!black,title={Before you move on}]
What does \bn{সাধ্য} mean? (\textquotedblleft{}One's power; what one can manage\textquotedblright{}.) Say it once more.
\end{tcolorbox}
\section[\texorpdfstring{\bn{মিতব্যয়ী} (mitôbyôyī)}{মিতব্যয়ী (mitôbyôyī)}]{\bn{মিতব্যয়ী} (mitôbyôyī) — thrifty}

\label{lesson:BN-C545-mitobyoyi}



\begin{quote}
Before the new one: say the Bengali for priority, then the Bengali for what one can manage.
\end{quote}

\subsection*{You'll want to know: \bn{মিতব্যয়ী}}
\textbf{\bn{মিতব্যয়ী}} — \emph{mitôbyôyī} — \textquotedblleft{}thrifty\textquotedblright{}.

\begin{quote}
\bn{মিতব্যয়ী} \bn{হলে} \bn{তুমি} \bn{টাকা} \bn{জমাতে} \bn{পারবে।}
\end{quote}

\begin{quote}
\emph{Mitôbyôyī hôle tumi ṭākā jômāte pārbe.}
\end{quote}

\begin{quote}
If you're thrifty, you'll be able to save money.
\end{quote}

\begin{cousinweb}
\textbf{\bn{মিতব্যয়ী}} (\emph{mitôbyôyī}) is a tatsama, Sanskrit \emph{mitavyayin}, \textquotedblleft{}spending little, frugal\textquotedblright{} (Monier-Williams): \emph{mita}, \textquotedblleft{}measured\textquotedblright{}, and \emph{vyaya}, \textquotedblleft{}spending\textquotedblright{}.
\end{cousinweb}

\begin{grammarlens}[title={What you've built}]
One more word for saying what happens if.
\end{grammarlens}

\subsection*{Your turn}
\begin{itemize}
\raggedright
  \item \emph{Say it:} \emph{mitôbyôyī}
  \item \emph{Say it:} \emph{mitôbyôyī}, once more
  \item \emph{Say it:} \emph{Mitôbyôyī hôle tumi ṭākā jômāte pārbe.} — If you're thrifty, you'll be able to save money.
  \item \emph{Recall:} say the Bengali for priority, then the Bengali for what one can manage, then say \emph{mitôbyôyī} again
\end{itemize}

\begin{tcolorbox}[breakable,colback=teal!4,colframe=teal!35!black,title={Before you move on}]
What does \bn{মিতব্যয়ী} mean? (\textquotedblleft{}Thrifty\textquotedblright{}.) Say it once more.
\end{tcolorbox}
\section[\texorpdfstring{\bn{ঠাওরানো} (ṭhāorāno)}{ঠাওরানো (ṭhāorāno)}]{\bn{ঠাওরানো} (ṭhāorāno) — to make out; to work out (what is right)}

\label{lesson:BN-C545-thaorano}



\begin{quote}
Before the new one: say the Bengali for what one can manage, then the Bengali for thrifty.
\end{quote}

\subsection*{You'll want to know: \bn{ঠাওরানো}}
\textbf{\bn{ঠাওরানো}} — \emph{ṭhāorāno} — \textquotedblleft{}to make out; to work out (what is right)\textquotedblright{}.

\begin{quote}
\bn{অন্ধকারে} \bn{আমি} \bn{রাস্তাটা} \bn{ঠাওরাতে} \bn{পারলাম} \bn{না।}
\end{quote}

\begin{quote}
\emph{Ôndhôkāre āmi rāstāṭā ṭhāorāte pārlām nā.}
\end{quote}

\begin{quote}
In the dark I couldn't make out the road.
\end{quote}

Past (he/she …): \bn{তিনি} \bn{রাস্তাটা} \bn{ঠাওরালেন}, \emph{tini rāstāṭā ṭhāorālen}.

\begin{cousinweb}
\textbf{\bn{ঠাওরানো}} (\emph{ṭhāorāno}) is a tadbhava: CDIAL lists Bengali \emph{ṭhāurāna}, \textquotedblleft{}to stop, to consider, to ascertain\textquotedblright{}, under Sanskrit \emph{sthāvara}, \textquotedblleft{}fixed, firm\textquotedblright{} (CDIAL 13767).
\end{cousinweb}

\begin{grammarlens}[title={What you've built}]
One more word for saying what happens if.
\end{grammarlens}

\subsection*{Your turn}
\begin{itemize}
\raggedright
  \item \emph{Say it:} \emph{ṭhāorāno}
  \item \emph{Say it:} \emph{ṭhāorāno}, once more
  \item \emph{Say it:} \emph{Ôndhôkāre āmi rāstāṭā ṭhāorāte pārlām nā.} — In the dark I couldn't make out the road.
  \item \emph{Recall:} say the Bengali for what one can manage, then the Bengali for thrifty, then say \emph{ṭhāorāno} again
\end{itemize}

\begin{tcolorbox}[breakable,colback=teal!4,colframe=teal!35!black,title={Before you move on}]
What does \bn{ঠাওরানো} mean? (\textquotedblleft{}To make out; to work out (what is right)\textquotedblright{}.) Say it once more.
\end{tcolorbox}
\section[(dialogue)]{Looking back: from Facing Obstacles to Choices}

\label{lesson:BN-R545-review-539-545}



\begin{quote}
Say the Bengali for a hindrance and persistent.
\end{quote}

\subsection*{Your turn}
Say these aloud:

\begin{itemize}
\raggedright
  \item a hindrance, persistent, desperate, to give up hope, to get tangled up
  \item a fee, the price of a mistake, reassurance, armour, to finish off
  \item expectation, long-lasting, reserved, a queue token, a phase
  \item to fall through, fruitless, in chaos, to come loose, a muddle
  \item accomplished, attainment, flawless, to secure, a desired goal
  \item prompt, a respite, to draw near, a scramble, in a jiffy
  \item one of several, priority, what one can manage, thrifty, to make out
\end{itemize}

\begin{tcolorbox}[breakable,colback=teal!4,colframe=teal!35!black,title={Before you move on}]
A hindrance? (\textbf{\bn{বিঘ্ন}}, \emph{bighnô}.) Persistent? (\textbf{\bn{নাছোড়বান্দা}}, \emph{nāchhoṛbāndā}.) Desperate? (\textbf{\bn{মরিয়া}}, \emph{môriyā}.) To give up hope? (\textbf{\bn{হাল} \bn{ছাড়া}}, \emph{hāl chhāṛā}.) To get tangled up? (\textbf{\bn{জড়ানো}}, \emph{jôṛāno}.) A fee? (\textbf{\bn{মাশুল}}, \emph{māshul}.) The price of a mistake? (\textbf{\bn{খেসারত}}, \emph{khesārôt}.) Reassurance? (\textbf{\bn{আশ্বাস}}, \emph{āshvās}.) Armour? (\textbf{\bn{বর্ম}}, \emph{bôrmô}.) To finish off? (\textbf{\bn{চুকানো}}, \emph{chukāno}.) Expectation? (\textbf{\bn{প্রতীক্ষা}}, \emph{prôtīkkhā}.) Long-lasting? (\textbf{\bn{দীর্ঘ}}, \emph{dīrghô}.) Reserved? (\textbf{\bn{সংরক্ষিত}}, \emph{sôngrôkkhit}.) A queue token? (\textbf{\bn{টোকেন}}, \emph{ṭoken}.) A phase? (\textbf{\bn{পর্যায়}}, \emph{pôrjāy}.) To fall through? (\textbf{\bn{ভেস্তে} \bn{যাওয়া}}, \emph{bheste jāoyā}.) Fruitless? (\textbf{\bn{নিষ্ফল}}, \emph{nishphôl}.) In chaos? (\textbf{\bn{বিশৃঙ্খল}}, \emph{bishringkhôl}.) To come loose? (\textbf{\bn{খসা}}, \emph{khôsā}.) A muddle? (\textbf{\bn{গোলমাল}}, \emph{golmāl}.) Accomplished? (\textbf{\bn{কৃতী}}, \emph{kritī}.) Attainment? (\textbf{\bn{সিদ্ধি}}, \emph{siddhi}.) Flawless? (\textbf{\bn{নিখুঁত}}, \emph{nikhũt}.) To secure? (\textbf{\bn{হাসিল} \bn{করা}}, \emph{hāsil kôrā}.) A desired goal? (\textbf{\bn{অভীষ্ট}}, \emph{ôbhīshṭô}.) Prompt? (\textbf{\bn{তৎপর}}, \emph{tôtpôr}.) A respite? (\textbf{\bn{অবকাশ}}, \emph{ôbôkāsh}.) To draw near? (\textbf{\bn{ঘনিয়ে} \bn{আসা}}, \emph{ghôniye āsā}.) A scramble? (\textbf{\bn{হুড়োহুড়ি}}, \emph{huṛohuṛi}.) In a jiffy? (\textbf{\bn{ঝটপট}}, \emph{jhôṭpôṭ}.) One of several? (\textbf{\bn{অন্যতম}}, \emph{ônyôtôm}.) Priority? (\textbf{\bn{অগ্রাধিকার}}, \emph{ôgrādhikār}.) What one can manage? (\textbf{\bn{সাধ্য}}, \emph{sādhyô}.) Thrifty? (\textbf{\bn{মিতব্যয়ী}}, \emph{mitôbyôyī}.) To make out? (\textbf{\bn{ঠাওরানো}}, \emph{ṭhāorāno}.)
\end{tcolorbox}
\section[(dialogue)]{Choices — once more}

\label{lesson:BN-R545-review-final-545}



\begin{quote}
Say the Bengali for one of several and priority.
\end{quote}

\subsection*{Your turn}
\begin{itemize}
\raggedright
  \item \emph{Say it:} one of several, priority, what one can manage, thrifty, to make out
\end{itemize}

\begin{tcolorbox}[breakable,colback=teal!4,colframe=teal!35!black,title={Before you move on}]
One of several? (\textbf{\bn{অন্যতম}}, \emph{ônyôtôm}.) Priority? (\textbf{\bn{অগ্রাধিকার}}, \emph{ôgrādhikār}.) What one can manage? (\textbf{\bn{সাধ্য}}, \emph{sādhyô}.) Thrifty? (\textbf{\bn{মিতব্যয়ী}}, \emph{mitôbyôyī}.) To make out? (\textbf{\bn{ঠাওরানো}}, \emph{ṭhāorāno}.)
\end{tcolorbox}
